% LỜI NÓI ĐẦU: đặt sau Lời cảm ơn, trước Mục lục; nêu lý do chọn đề tài và cách đọc báo cáo, không lặp lại Chương 1.
\clearpage
\section*{LỜI NÓI ĐẦU}
\addcontentsline{toc}{section}{Lời nói đầu}

Chuỗi thời gian xuất hiện ở hầu hết các lĩnh vực, từ mực nước hồ và lưu lượng sông đến giá chứng khoán và lượng hành khách hàng không. Các mô hình ARMA và ARIMA là công cụ cổ điển để mô tả và dự báo những chuỗi này, nhưng lý thuyết của chúng dựa trên các giả thiết mà dữ liệu thực không phải lúc nào cũng thỏa mãn. Vì vậy nhóm chọn cách học bằng thực nghiệm: lấy từng kết quả lý thuyết, chạy nó trên dữ liệu thực bằng Python, rồi đối chiếu xem lý thuyết nói gì và dữ liệu cho thấy gì.

Báo cáo này là kết quả của quá trình đó. Thay vì nhắc lại toàn bộ lý thuyết, báo cáo chỉ tóm tắt những gì cần cho việc kiểm chứng ở \chg{2}, rồi đi thẳng vào kiểm chứng ở \chg{4}, \chn{5} và \chn{6}. Mỗi kết quả đều có đủ phần phát biểu, cách làm, mã lệnh, số liệu, diễn giải và nhận xét, nên người đọc có thể theo dõi hoặc tự chạy lại. Nơi nào kết quả khớp lý thuyết, nhóm nói rõ; nơi nào lệch, nhóm cố gắng chỉ ra lý do.

Để đọc thuận tiện, các bộ dữ liệu được gọi bằng những ký hiệu ngắn như HU, LY hay DR, và mọi chữ viết tắt đều được giải thích trong Danh mục ký hiệu và chữ viết tắt ngay sau Mục lục. Người đã quen với lý thuyết có thể đọc thẳng các chương kiểm chứng, còn các nhận xét tổng hợp nằm ở \chg{7} và \chg{8}. Mã nguồn đầy đủ và cách chạy lại có trong \phuluc{B}.

Nhóm hy vọng báo cáo này vừa chính xác vừa dễ theo dõi, và mong nhận được những góp ý để hoàn thiện hơn.
